In this paper, we introduced the central distribution as a representative probability distribution for a set or
hierarchy of admissible distributions.
The central distribution is characterized by a minimax-KL principle, providing a principled probabilistic
representation of partial knowledge while remaining within the standard framework of Bayesian probability theory.

For finite sets of distributions, we established existence, uniqueness, and characterization results and showed that
central distributions admit a convex-combination representation.
We extended the framework to hierarchies of distributions and introduced the notion of a central prior for
parametric families.
To address problems involving growing sequences of observations, we proposed the concept of asymptotic central
distributions and derived sufficient conditions for asymptotic centrality.
The resulting asymptotic framework suggests a connection between central priors and classical noninformative priors,
including Jeffreys' prior in several standard models.
In addition, we developed approximation bounds that lead to practical procedures for constructing approximate
central distributions.

The framework suggests a new way of representing uncertainty about probabilistic models themselves.
Rather than selecting a single model or beginning with a prior distribution over models, uncertainty may be
represented directly through sets or hierarchies of admissible distributions and summarized by their corresponding
central distributions.
From this perspective, central distributions provide a bridge between robust, information-theoretic, and Bayesian
approaches.

Several important questions remain open.
In particular, a rigorous extension of the theory to general parametric families,
computational methods for computing and approximating central distributions,
and a rigorous asymptotic theory for central priors require further investigation.
It would also be valuable to study central distributions in a broader range of practical statistical problems in
order to better understand their computational and inferential properties.
We hope that the central distribution framework provides a useful foundation for studying noninformative priors and
hierarchical partial knowledge from an information-theoretic perspective.
